\section{Постановка задачи}\label{sec2}

\subsection{Общая постановка мультимодального декодирования}\label{sec2:general}

Пусть $\mathcal{X}_{(m)}$~--- пространство наблюдений модальности $m$, где $m = 1, \ldots, M$. Совместное пространство нейробиологических данных определяется как декартово произведение:
\begin{equation}\label{eq:joint_space}
\mathcal{X} = \mathcal{X}_{(1)} \times \cdots \times \mathcal{X}_{(M)}.
\end{equation}
Пусть $\mathcal{Y}$~--- пространство визуальных стимулов. Задача мультимодального декодирования состоит в построении отображения
\begin{equation}\label{eq:decoding}
\mathbf{g}: \mathcal{X} \to \mathcal{Y},
\end{equation}
которое восстанавливает стимул $\mathbf{y} \in \mathcal{Y}$ по совместному наблюдению $\mathbf{x} \in \mathcal{X}$. Настоящая работа рассматривает частный случай $M = 2$, в котором модальностями являются одновременные записи фМРТ и ЭЭГ.

\subsection{Двумодальное декодирование фМРТ--ЭЭГ}\label{sec2:bimodal}

Задан датасет триплетов (фМРТ, ЭЭГ, изображение):
\begin{equation}\label{eq:dataset}
    \mathfrak{D} = \{(\mathbf{F}_i, \mathbf{E}_i, \mathbf{I}_i) \mid i = 1, \ldots, N \}, \quad (\mathbf{F}_i, \mathbf{E}_i) \in \mathcal{X}_f \times \mathcal{X}_e, \quad \mathbf{I}_i \in \mathcal{Y},
\end{equation}
где $\mathbf{F}_i \in \mathbb{R}^{X \times Y \times Z}$~--- данные фМРТ с пространственными размерностями $X$, $Y$ и $Z$; $\mathbf{E}_i \in \mathbb{R}^{K \times T}$~--- данные ЭЭГ, где $K$~--- число каналов, $T$~--- длина временного окна; $\mathbf{I}_i \in \mathbb{R}^{H \times W \times C}$~--- изображение с высотой $H$, шириной $W$ и количеством каналов $C$. Требуется построить отображение $\mathbf{f}$ из совместного пространства данных фМРТ--ЭЭГ в пространство изображений:
\begin{equation}\label{eq:target}
    \mathbf{f}:  \mathcal{X}_f \times \mathcal{X}_e \rightarrow \mathcal{Y}.
\end{equation}

\subsection{Гипотеза о гемодинамической задержке}\label{sec2:dt}

BOLD-сигнал фМРТ возникает с гемодинамической задержкой $\Delta t$ относительно момента предъявления визуального стимула~\cite{Bandettini1992, Logothetis2003}; ЭЭГ-сигнал, напротив, отражает нейронный отклик практически без задержки. При формировании триплета $(\mathbf{F}_i, \mathbf{E}_i, \mathbf{I}_i)$ временной сдвиг затрагивает только фМРТ-составляющую: изображение $\mathbf{I}_i$ должно соответствовать моменту времени, предшествующему моменту регистрации фМРТ-скана $\mathbf{F}_i$ на величину $\Delta t$. Формально пусть $\mathbf{F}(t)$~--- фМРТ-скан в момент $t$, $\mathbf{E}(t)$~--- ЭЭГ-окно вокруг момента $t$, а $\mathbf{I}(t)$~--- изображение, предъявленное в момент $t$. Тогда триплеты выборки~\eqref{eq:dataset} формируются по правилу
\begin{equation}\label{eq:shift}
\mathbf{F}_i = \mathbf{F}(t_i), \qquad \mathbf{E}_i = \mathbf{E}(t_i - \Delta t), \qquad \mathbf{I}_i = \mathbf{I}(t_i - \Delta t).
\end{equation}
В качестве значения $\Delta t$ принимается оценка $4$--$6$~с, согласующаяся с известным медицинским диапазоном гемодинамической задержки~\cite{Bandettini1992} и эмпирически подтвержденная в нашей предыдущей работе~\cite{dorin2024forecasting}.
